\chapter{Kryteria planarności, podziały krawędzi i kolorowanie}
\label{ch:tg-kuratowski}

\section{Podgrafy, podziały krawędzi i minory}
\label{sec:tg-rozne-redukcje}

\begin{definicja}[Podział krawędzi i podgraf topologiczny]
Niech $G$ i $H$ będą skończonymi grafami prostymi.
\emph{Podziałem krawędzi}
$uv\in E(H)$ nazywamy zastąpienie tej krawędzi ścieżką
$u-w-v$, gdzie $w$ jest nowym wierzchołkiem stopnia dwa.
Graf $H'$ jest \emph{subdywizją} $H$, jeżeli jest izomorficzny
z grafem otrzymanym z $H$ przez skończony ciąg takich
podziałów (dopuszczamy ciąg pusty). Mówimy, że $H$ jest \emph{podgrafem
topologicznym} grafu $G$, jeżeli pewna subdywizja $H$
jest izomorficzna z podgrafem $G$.
\end{definicja}

\begin{stwierdzenie}[Niezmienniczość planarności przy subdywizji]
Niech $H$ będzie skończonym grafem prostym, a $H'$
grafem otrzymanym z $H$ przez podział jednej krawędzi. Wtedy $H$ jest planarny wtedy i tylko wtedy,
gdy $H'$ jest planarny.
\end{stwierdzenie}

\begin{proof}
Jeśli znamy zanurzenie $H$, wybieramy dowolny punkt we wnętrzu
obrazu dzielonej krawędzi i traktujemy go jako nowy wierzchołek.
Tak powstaje zanurzenie $H'$.
W drugą stronę, w zanurzeniu $H'$ nowy wierzchołek ma stopień
dwa. Usuwamy jego oznaczenie i traktujemy dwie incydentne
krawędzie jako jeden łuk prosty: ich obrazy mają wspólny
tylko usuwany wierzchołek, a pozostałe końce są różne.
Nie pojawiają się nowe przecięcia, więc otrzymujemy zanurzenie $H$.
\end{proof}

\begin{definicja}[Minor grafu]
Niech $G$ i $H$ będą skończonymi grafami. Mówimy, że
$H$ jest \emph{minorem} $G$, jeśli jest izomorficzny
z grafem otrzymanym z $G$ przez skończony ciąg operacji
usuwania wierzchołków wraz z ich incydentnymi krawędziami,
usuwania krawędzi i \emph{kontrakcji} krawędzi o różnych
końcach (utożsamienia obu jej końców po usunięciu tej
krawędzi). Dopuszczamy ciąg pusty. W grafach pośrednich
mogą wystąpić pętle i krawędzie równoległe, które można
następnie usunąć. Jeżeli kontrakcja dotyczy grafu prostego
i ma dać graf prosty, po utożsamieniu końców usuwamy
pętle i pozostawiamy po jednej krawędzi z każdej klasy
krawędzi równoległych.
\end{definicja}

\begin{stwierdzenie}[Planarność jest zamknięta na minory]
Jeżeli skończony graf $G$ jest planarny i $H$ jest jego minorem,
to $H$ jest planarny.
\end{stwierdzenie}

\begin{proof}
Usuwanie krawędzi i wierzchołków nie wprowadza przecięć.
Dla kontraktowanej krawędzi o różnych końcach bierzemy
dysk złożony z dwóch dysków wierzchołkowych i pasa
tej krawędzi, korzystając z regularnego otoczenia
opisanego w dowodzie twierdzenia \ref{thm:tg-euler-plane}.
Pozostałe incydencje przecinają jego brzeg w różnych
punktach. Wewnątrz dysku zastępujemy dotychczasowy
rysunek jednym nowym wierzchołkiem połączonym
promieniami z tymi punktami, zachowując rysunek poza
dyskiem. Jeżeli inna krawędź ma oba końce w dysku, jej
fragment poza dyskiem łączy dwa różne punkty brzegu,
a po operacji krawędź staje się pętlą. Wszystkie
nowe odcinki wewnątrz dysku mają wspólny wyłącznie
nowy wierzchołek. Otrzymujemy osadzenie grafu po
kontrakcji; usunięcie pętli lub nadmiarowych krawędzi
równoległych również nie tworzy przecięć.
Stąd każda operacja minorowa zachowuje planarność.
\end{proof}

\begin{twierdzenie}[Kuratowski]
\label{thm:tg-kuratowski}
Niech $G$ będzie skończonym grafem prostym. Następujące
warunki są równoważne:
\begin{enumerate}[label=(\roman*)]
\item $G$ jest planarny;
\item $G$ nie zawiera podgrafu izomorficznego z żadną
subdywizją $K_5$ ani $K_{3,3}$.
\end{enumerate}
\end{twierdzenie}

\begin{proof}[Dowód implikacji \texorpdfstring{$(i)\Rightarrow(ii)$}{(i) => (ii)}]
Gdyby planarny $G$ zawierał podgraf będący subdywizją
$K_5$ lub $K_{3,3}$, sam ten podgraf byłby planarny
(jako podgraf grafu planarnego). Cofając podziały krawędzi,
uzyskalibyśmy zanurzenie $K_5$ albo $K_{3,3}$.
Sprzeciwia się to obliczeniu z poprzedniego rozdziału,
w którym wykazaliśmy ich nieplanarność.
\end{proof}

\section{Pełny dowód twierdzenia Kuratowskiego}
\label{sec:tg-kuratowski-full-proof}

Przedstawimy samodzielny dowód implikacji
$(ii)\Rightarrow(i)$ z twierdzenia \ref{thm:tg-kuratowski}.
Oparty jest na kontrakcji krawędzi zachowującej
trójspójność i analizie osadzenia grafu po kontrakcji.
Wymaga jawnego odróżnienia \emph{minora} od
\emph{subdywizji}; odpowiedni pomost udowodnimy najpierw.

\begin{definicja}[Model minora]
\label{def:tg-kuratowski-minor-model}
Niech $G,H$ będą skończonymi grafami prostymi.
\emph{Modelem} $H$ w $G$ nazywamy rodzinę
$(W_h)_{h\in V(H)}$ parami rozłącznych,
niepustych podzbiorów $V(G)$, takich że każdy
podgraf indukowany $G[W_h]$ jest spójny, a ponadto
dla każdej krawędzi $hh'\in E(H)$ istnieje
krawędź $G$ o jednym końcu w $W_h$,
a drugim w $W_{h'}$.
Zbiory $W_h$ nazywamy \emph{zbiorami rozgałęzienia}.
\end{definicja}

\begin{lemat}[Model a minor]
\label{lem:tg-kuratowski-model-minor}
Dla skończonych grafów prostych $G,H$
graf $H$ jest minorem $G$ wtedy i tylko wtedy,
gdy istnieje model $H$ w $G$.
\end{lemat}
\begin{proof}
Mając model, usuwamy najpierw wszystkie wierzchołki
spoza $\bigcup_{h\in V(H)}W_h$. W każdym spójnym
$G[W_h]$ wybieramy drzewo rozpinające i kontraktujemy
jego krawędzie. Rozłączność zbiorów $W_h$ gwarantuje,
że każdy z nich daje dokładnie jeden wierzchołek
i żadne dwa różne zbiory nie zostają utożsamione.
Każdej krawędzi $hh'\in E(H)$ odpowiada co najmniej
jedna krawędź między uzyskanymi wierzchołkami.
Zachowujemy po jednej takiej krawędzi, usuwając
wszystkie pozostałe, w tym pętle i nadmiarowe
krawędzie równoległe. Otrzymujemy graf izomorficzny z $H$.

Odwrotnie, rozpatrzmy ciąg operacji
prowadzących z $G$ do $H$.
Każdemu wierzchołkowi grafu pośredniego
przypisujemy zbiór jego pierwotnych
wierzchołków: początkowo są to singletony,
zaś kontrakcja krawędzi łączy dwa
takie zbiory. Utrzymujemy dwa niezmienniki: każdy zbiór indukuje
spójny podgraf w \emph{pierwotnym} grafie $G$, a każda
krawędź między różnymi wierzchołkami grafu pośredniego
pochodzi od krawędzi $G$ między ich zbiorami.
Początkowo oba warunki są oczywiste. Przy kontrakcji
drugi warunek zapewnia krawędź $G$ łączącą dwa zbiory,
więc ich suma indukuje spójny podgraf $G$.
Każda zachowana krawędź między różnymi wierzchołkami
zachowuje swojego świadka w $G$; przy zastąpieniu
krawędzi równoległych jedną zachowujemy jednego
z ich świadków. Usunięcie wierzchołka usuwa jego
zbiór, a usunięcie krawędzi nie zmienia zachowanych
zbiorów ani świadków. Spójność rozważamy zawsze
w $G$, a nie w grafie pośrednim. Końcowe zbiory
są parami rozłączne i odpowiadają
wierzchołkom $H$; dla każdej
krawędzi $H$ istnieje pierwotna
krawędź między właściwymi zbiorami.
Są więc modelem $H$ w $G$.
\end{proof}

\begin{lemat}[Trzy lub cztery zaznaczone końce drzewa]
\label{lem:tg-kuratowski-tree-terminals}
Niech $T$ będzie skończonym drzewem,
a na jego wierzchołkach umieszczono
$k\in\{3,4\}$ różnych etykiet;
kilka etykiet może leżeć
na tym samym wierzchołku.
Wówczas zachodzi dokładnie jedna
z następujących możliwości:
\begin{enumerate}[label=(\roman*)]
\item istnieje wierzchołek $z\in V(T)$
taki, że ścieżki od $z$ do wierzchołków
zaznaczonych etykietami mają parami część wspólną
równą $\{z\}$. Ścieżka do etykiety leżącej w $z$
jest ścieżką długości zero; dopuszczamy więc kilka
takich ścieżek;
\item $k=4$ i pewna krawędź drzewa
rozdziela etykiety na dwie pary.
\end{enumerate}
\end{lemat}
\begin{proof}
Przy $k=3$ usunięcie dowolnej
krawędzi dzieli drzewo na dwie
części, z których co najmniej
jedna zawiera co najmniej
dwie etykiety. Orientujemy
krawędź ku tej części.
Przy $k=4$ albo już
występuje podział $2+2$
i mamy (ii), albo żadna
krawędź nie ma takiego
podziału; wówczas
kierujemy każdą krawędź
ku stronie zawierającej
co najmniej trzy etykiety.
W obu przypadkach
orientacja każdej krawędzi
jest określona jednoznacznie.

W skończonym skierowanym
drzewie istnieje wierzchołek
$z$ bez wychodzących krawędzi:
idąc kolejno po krawędziach
zgodnie z orientacją,
nie możemy powrócić do
wcześniejszego wierzchołka,
ponieważ drzewo nie ma cykli.
Dla każdej składowej drzewa
$T-z$ krawędź łącząca
ją z $z$ jest skierowana ku
$z$, a zatem ta składowa
zawiera co najwyżej jedną
etykietę. Ścieżki od $z$
do zaznaczonych wierzchołków
biegną więc w różnych
składowych $T-z$
(poza etykietami znajdującymi
się w samym $z$).
To daje (i).
Jeśli pewna krawędź rozdziela cztery etykiety na dwie
pary, dowolny kandydat na $z$ leży w jednej z dwóch
składowych po usunięciu tej krawędzi. Obie ścieżki
od $z$ do etykiet po drugiej stronie muszą zawierać
tę krawędź, a więc również jej koniec różny od $z$.
Nie spełniają warunku (i). Dla trzech etykiet
przypadek (ii) jest wykluczony przez jego sformułowanie.
Obie możliwości są zatem rozłączne.
\end{proof}

\begin{lemat}[Od zabronionych minorów do subdywizji]
\label{lem:tg-kuratowski-minor-to-subdivision}
Niech $G$ będzie skończonym grafem prostym.
\begin{enumerate}[label=(\roman*)]
\item Jeśli $K_{3,3}$ jest minorem $G$,
to $G$ zawiera podgraf izomorficzny
z subdywizją $K_{3,3}$.
\item Jeśli $K_5$ jest minorem $G$,
to $G$ zawiera podgraf izomorficzny
z subdywizją $K_5$ albo $K_{3,3}$.
\end{enumerate}
\end{lemat}
\begin{proof}
Niech $H$ oznacza jeden z grafów
$K_{3,3},K_5$. Wybierzmy model
$(W_h)_{h\in V(H)}$ w $G$
z lematu \ref{lem:tg-kuratowski-model-minor}.
Dla każdej nieuporządkowanej krawędzi $hh'\in E(H)$
wybierzmy jedną krawędź $e_{hh'}=e_{h'h}$ grafu $G$
łączącą $W_h$ z $W_{h'}$. Jej końce oznaczmy
odpowiednio przez $a_{h,h'}\in W_h$ oraz
$a_{h',h}\in W_{h'}$. Różnym krawędziom $H$
odpowiadają różne krawędzie $G$, gdyż ich końce
należą do różnych par zbiorów rozgałęzienia.
W spójnym grafie $G[W_h]$
wybierzmy drzewo rozpinające $T_h$
i zaznaczmy na nim etykietę $h'$
w wierzchołku $a_{h,h'}$,
dla każdego sąsiada $h'$ w $H$.
Każde $T_h$ ma zatem dokładnie
$\deg_H(h)$ etykiet, choć
niekoniecznie na różnych
wierzchołkach.

Jeżeli dla każdego drzewa $T_h$
zachodzi przypadek (i) lematu
\ref{lem:tg-kuratowski-tree-terminals},
oznaczmy uzyskany środek
przez $z_h$. Dla każdej
krawędzi $hh'$ łączymy
$z_h$ z $z_{h'}$ ścieżką
złożoną z odcinka w $T_h$
do $a_{h,h'}$, wybranej
krawędzi między zbiorami
rozgałęzienia oraz odcinka
w $T_{h'}$ do $z_{h'}$.
Każda tak skonstruowana ścieżka $P_{hh'}$ jest
prosta: jej części drzewowe są proste i leżą
w rozłącznych zbiorach $W_h,W_{h'}$, a łączy je
jedna krawędź $e_{hh'}$. Jeżeli któraś etykieta
leży w środku drzewa, odpowiednia część drzewowa
ma długość zero; krawędź $e_{hh'}$ nadal zapewnia
dodatnią długość całej ścieżki.
Dla dwóch różnych krawędzi $H$ bez wspólnego końca
ich ścieżki są rozłączne, ponieważ używają
rozłącznych zbiorów rozgałęzienia. Gdy krawędzie
$H$ mają wspólny koniec $h$, ścieżki spotykają się
dokładnie w $z_h$ na mocy warunku (i) lematu
\ref{lem:tg-kuratowski-tree-terminals}.
Ponadto wnętrze $P_{hh'}$ nie zawiera żadnego
z wierzchołków $z_j$: dla $j=h,h'$ wynika to
z prostoty ścieżki, a dla pozostałych $j$
z rozłączności zbiorów rozgałęzienia.
Bierzemy podgraf o zbiorze krawędzi równym sumie
zbiorów krawędzi wszystkich tych ścieżek i o zbiorze
wierzchołków równym ich sumie wierzchołkowej.
Nie bierzemy dodatkowych krawędzi grafu $G$
między tymi wierzchołkami. Każda ścieżka zastępuje
jedną krawędź $H$ przez ciąg jej podziałów, więc
otrzymany podgraf jest izomorficzny z subdywizją $H$.

Dla $H=K_{3,3}$ wszystkie
wierzchołki mają stopień
trzy, więc dla każdego
$T_h$ lemat zapewnia (i).
Dowodzi to pierwszej tezy.

Niech teraz $H=K_5$.
Jeśli wszystkie drzewa
spełniają (i), uzyskujemy
subdywizję $K_5$.
W przeciwnym razie
pewne drzewo, powiedzmy
$T_1$, ma krawędź rozcinającą jego cztery etykiety
na dwie pary. Ponieważ każda permutacja wierzchołków
$K_5$ jest automorfizmem, po zmianie numeracji
pozostałych czterech wierzchołków możemy przyjąć,
że są to pary $\{2,3\}$ i $\{4,5\}$.
Niech $U$ i $W$ będą
zbiorami wierzchołków
dwóch składowych po
usunięciu tej krawędzi,
z etykietami odpowiednio
$2,3$ oraz $4,5$.
Poniższe rodziny po
trzy zbiory
\[
\{\,U,W_4,W_5\,\},
\qquad
\{\,W,W_2,W_3\,\}
\]
są dwiema częściami
modelu $K_{3,3}$:
jako świadków dziewięciu wymaganych krawędzi
bierzemy krawędź rozcinającą $T_1$, krawędzie
$e_{12},e_{13}$ o końcach w $U$, krawędzie
$e_{14},e_{15}$ o końcach w $W$ oraz krawędzie
$e_{24},e_{25},e_{34},e_{35}$.
Zbiory $U,W$ są niepuste, rozłączne i spójne
w $G$, gdyż indukowane przez nie podgrafy $G$
zawierają odpowiednie składowe drzewa $T_1$.
Pozostałe cztery zbiory są niepuste i spójne
z definicji modelu, a wszystkie sześć zbiorów
jest parami rozłącznych. Nie wymaga się braku
krawędzi pomiędzy zbiorami z tej samej części:
takie krawędzie nie są potrzebne w modelu minora.
Zatem $G$ ma minor
$K_{3,3}$, który na mocy
już udowodnionego
przypadku (i) prowadzi
do jego subdywizji.
\end{proof}

\begin{wniosek}[Przeszkody są dziedziczone przez minory]
\label{cor:tg-kuratowski-minor-heredity}
Jeżeli skończony graf prosty $G$
nie zawiera subdywizji $K_5$
ani $K_{3,3}$ jako podgrafu,
to żaden minor $G$
nie zawiera takiej subdywizji.
\end{wniosek}
\begin{proof}
Gdyby minor $G'$ zawierał
którąś subdywizję,
to miałby $K_5$ lub $K_{3,3}$ jako minor:
usuwamy wszystkie krawędzie i wierzchołki spoza
subdywizji, a następnie na każdej ścieżce
zastępującej krawędź kontraktujemy wszystkie
krawędzie oprócz jednej. Rozłączność wnętrz
ścieżek pozwala wykonać te kontrakcje bez
utożsamiania różnych wierzchołków rozgałęzienia.
Składając ciąg operacji prowadzących od $G$ do $G'$
z tym ciągiem, otrzymujemy odpowiedni minor w $G$.
Lemat
\ref{lem:tg-kuratowski-minor-to-subdivision}
prowadziłby do zabronionej
subdywizji w $G$.
\end{proof}


\begin{definicja}[Trójspójność]
\label{def:tg-kuratowski-threeconnected}
Skończony graf prosty $G$ jest
\emph{trójspójny wierzchołkowo},
jeśli $|V(G)|\geq4$, jest spójny
i usunięcie dowolnego zbioru
co najwyżej dwóch wierzchołków
nie rozspójnia pozostałego grafu.
W dalszych lematach zapis $G/xy$ oznacza prosty
graf otrzymany przez kontrakcję krawędzi $xy$,
z usunięciem pętli i nadmiarowych krawędzi
równoległych, zgodnie z powyższą konwencją.
\end{definicja}

\begin{lemat}[Małe separatory grafu trójspójnego]
\label{lem:tg-kuratowski-threecut}
Niech $G$ będzie skończonym
grafem trójspójnym,
a $S\subseteq V(G)$
zbiorem trzech wierzchołków.
Jeśli $G-S$ ma co najmniej
dwie składowe, to każdy
wierzchołek $s\in S$ ma
sąsiada w każdej
składowej $G-S$.
\end{lemat}
\begin{proof}
Wybierzmy składową $C$ grafu
$G-S$ i $s\in S$.
Gdyby $s$ nie miał
sąsiada w $C$,
to usunięcie $S\setminus\{s\}$
pozostawiałoby wszystkie
wierzchołki $C$ odłączone
od $s$ i pozostałych
składowych $G-S$.
Otrzymalibyśmy separator
złożony z dwóch wierzchołków,
co przeczy trójspójności.
\end{proof}

\begin{lemat}[Kontrakcja zachowująca trójspójność]
\label{lem:tg-kuratowski-contractible-edge}
Jeżeli $G$ jest skończonym
grafem prostym trójspójnym
i $G\not\cong K_4$,
to istnieje krawędź $uv$,
dla której po kontrakcji
graf $G/uv$ jest trójspójny.
\end{lemat}
\begin{proof}
Trójspójny graf o czterech wierzchołkach jest
izomorficzny z $K_4$: gdyby dwa wierzchołki nie
były sąsiednie, usunięcie pozostałych dwóch
pozostawiałoby graf niespójny. Zatem $|V(G)|\geq5$.

\emph{Krok 1: kryterium kontraktowalności.}
Dla dowolnej krawędzi $xy$ oznaczmy przez $v_{xy}$
wierzchołek powstały po kontrakcji. Wykażemy równoważność
\[
G/xy\text{ jest trójspójny}
\quad\Longleftrightarrow\quad
G-\{x,y,z\}\text{ jest spójny dla każdego }
z\in V(G)\setminus\{x,y\}.
\]
Graf $G/xy$ jest spójny i ma co najmniej cztery
wierzchołki. Jeżeli $S\subseteq V(G/xy)$,
$|S|\leq2$ i $v_{xy}\notin S$, utożsamiamy $S$
z tym samym zbiorem w $G$. Wtedy $(G/xy)-S$
jest wynikiem kontrakcji krawędzi $xy$ w spójnym
grafie $G-S$, a więc jest spójny.
Jeżeli $S=\{v_{xy}\}$, mamy izomorfizm
$(G/xy)-S\cong G-\{x,y\}$, więc również w tym
przypadku pozostały graf jest spójny.
Jedynym nierozstrzygniętym przypadkiem jest
$S=\{v_{xy},z\}$, gdzie $z\notin\{x,y\}$.
Wtedy $(G/xy)-S\cong G-\{x,y,z\}$: pozostałe
wierzchołki i krawędzie między nimi nie zostały
zmienione przez kontrakcję ani uproszczenie grafu.
To dowodzi równoważności.

\emph{Krok 2: minimalna składowa.}
Przypuśćmy przeciwnie, że żadna kontrakcja
nie zachowuje trójspójności. Na mocy kroku 1
dla każdej krawędzi $xy$ istnieje $z\notin\{x,y\}$,
dla którego $G-\{x,y,z\}$ ma co najmniej dwie
składowe. Wśród wszystkich trójek $(xy,z,C)$,
gdzie $xy\in E(G)$, $z\notin\{x,y\}$,
$G-\{x,y,z\}$ jest niespójny, a $C$ jest jedną
z jego składowych, wybierzmy trójkę minimalizującą
$|V(C)|$. Taki wybór istnieje, gdyż zbiór trójek
jest skończony i niepusty.
Piszemy dalej $C$ także na oznaczenie zbioru jego
wierzchołków. Lemat
\ref{lem:tg-kuratowski-threecut} zapewnia sąsiada
$w\in C$ wierzchołka $z$.

\emph{Krok 3: ścisłe zmniejszenie.}
Z założenia kontrakcja krawędzi $zw$ również
nie zachowuje trójspójności. Z kroku 1 istnieje
$t\notin\{z,w\}$, dla którego $G-\{z,w,t\}$
ma co najmniej dwie składowe.
Jeśli $t\notin\{x,y\}$, oba wierzchołki $x,y$
pozostają w tym grafie i należą do tej samej
składowej, gdyż łączy je krawędź $xy$.
Jeśli $t=x$ albo $t=y$, pozostaje dokładnie
jeden z nich. W obu przypadkach co najmniej
jeden z $x,y$ pozostaje, ponieważ $z,w$ są od
nich różne. Możemy zatem wybrać składową $D$
grafu $G-\{z,w,t\}$ niezawierającą żadnego
z wierzchołków $x,y$.

Z lematu \ref{lem:tg-kuratowski-threecut}
wierzchołek $w$ ma sąsiada $d\in D$.
Mamy $w\in C$ i $d\notin\{x,y,z\}$, więc
krawędź $wd$ leży w $G-\{x,y,z\}$ i $d\in C$.
Każdy wierzchołek $D$ można połączyć z $d$
ścieżką w $D$. Ponieważ $D$ nie zawiera
$x,y,z$, cała ta ścieżka leży w $G-\{x,y,z\}$,
a zatem $D\subseteq C$. Wierzchołek $w$ należy
do $C$, ale nie do $D$, gdyż został usunięty.
Stąd $0<|D|<|C|$. Argument obejmuje zarówno $t\in C$, jak i
$t\in\{x,y\}$. Inne położenie $t$ jest w istocie
niemożliwe: z lematu
\ref{lem:tg-kuratowski-threecut} również $t$ ma
sąsiada w $D\subseteq C$. Gdyby $t$ nie należał
ani do $C$, ani do $\{x,y,z\}$, byłby w innej
składowej $G-\{x,y,z\}$, co wyklucza taką krawędź.
Ponieważ $t\ne z$, wynika stąd
$t\in C\cup\{x,y\}$.

Trójka $(zw,t,D)$ należy do tej samej rodziny,
nad którą minimalizowaliśmy w kroku 2.
Otrzymana ścisła nierówność przeczy minimalności
$|C|$. Zatem istnieje krawędź, której kontrakcja
zachowuje trójspójność.
\end{proof}

\begin{lemat}[Sklejanie dwóch osadzeń]
\label{lem:tg-kuratowski-gluing}
Niech $G=G_1\cup G_2$ będzie
skończonym grafem prostym i niech
$V(G_1)\cap V(G_2)$
zawiera co najwyżej dwa
wierzchołki. Zakładamy,
że nie ma krawędzi
łączącej wierzchołek
należący wyłącznie
do $G_1$ z wierzchołkiem
należącym wyłącznie
do $G_2$.
\begin{enumerate}[label=(\roman*)]
\item Jeśli przecięcie ma
co najwyżej jeden
wierzchołek i oba $G_i$
są planarne, to $G$
jest planarny.
\item Jeśli przecięcie
składa się z dwóch
różnych wierzchołków
$x,y$ oraz oba grafy
$G_1+xy$ i $G_2+xy$
są planarne (przez
$+xy$ rozumiemy
dołączenie krawędzi,
o ile jeszcze jej nie ma),
to $G$ jest planarny.
\end{enumerate}
\end{lemat}
\begin{proof}
Przypadek pustego przecięcia
realizujemy przez umieszczenie
dwóch rysunków w rozłącznych
dyskach. Dla wspólnego wierzchołka $v$ rozpatrujemy oba
osadzenia na osobnych sferach. W każdym wybieramy
sektor regularnego otoczenia przy $v$ wolny od krawędzi
(i dowolny sektor, jeśli $v$ jest izolowany).
W tym sektorze umieszczamy mały domknięty dysk,
którego jedynym punktem wspólnym z grafem jest $v$
na brzegu dysku. Domknięcie dopełnienia tego dysku
w sferze jest, na mocy twierdzenia Jordana--Schoenfliesa,
dyskiem zawierającym cały graf i mającym $v$ na brzegu.
Istnienie użytego sektora wynika z lokalnego prostowania
grafu w powierzchni, przyjętego w dowodzie twierdzenia
\ref{thm:tg-euler-plane}.
Przenosimy te dwa dyski homeomorfizmami na dwa
domknięte dyski w płaszczyźnie styczne w jednym
punkcie, tak aby oba obrazy $v$ były tym punktem.
Poza nim dyski są rozłączne, więc identyfikacja
wystąpień $v$ nie tworzy żadnego nowego przecięcia.

Jeżeli grafy mają wspólne
końce $x,y$, rysujemy
najpierw $G_1+xy$
i $G_2+xy$ na sferze.
W każdym osadzeniu wybieramy regularne otoczenie
krawędzi $xy$ z dyskami przy jej końcach i pasem
wzdłuż jej wnętrza, jak w dowodzie twierdzenia
\ref{thm:tg-euler-plane}. Po jednej stronie pasa
prowadzimy drugi łuk od $x$ do $y$, wykorzystując
przy końcach sektory bezpośrednio przylegające do
krawędzi $xy$. Wnętrze łuku jest rozłączne z całym
grafem, także wtedy, gdy $xy$ jest mostem. Oba łuki ograniczają
mały dysk o wnętrzu rozłącznym z rysunkiem.
Domknięcie drugiej strony tej krzywej Jordana
jest dyskiem zawierającym cały graf, którego
brzeg zawiera $xy$ i nie zawiera żadnego innego
fragmentu grafu. Każdy z tych dysków przenosimy na jedną z dwóch
połówek standardowego domkniętego dysku, tak by
krawędź $xy$ została odwzorowana na wspólną średnicę,
a $x,y$ na jej ustalone końce. Można to zrobić,
wybierając najpierw zgodne homeomorfizmy brzegów
i przedłużając je promieniowo na dyski (po identyfikacji
każdego z nich z dyskiem standardowym).
Obrazy obu rysunków mają więc wspólną dokładnie
krawędź $xy$ z jej końcami.
Wszystkie pozostałe łuki
leżą w rozłącznych
obszarach, więc nie mają
nowych przecięć.
Jeśli $xy$ nie było
krawędzią $G$,
usuwamy pomocniczy łuk;
jeśli było, pozostawiamy
tylko jedno jego
wystąpienie.
Otrzymujemy osadzenie $G$.
Konstrukcja korzysta
ze zwykłego własnego
otoczenia łuku w sferze
i twierdzenia Jordana.
\end{proof}

\begin{lemat}[Ściana grafu dwuspójnego]
\label{lem:tg-kuratowski-face-cycle}
Niech $H$ będzie skończonym prostym
grafem dwuspójnym o co
najmniej trzech wierzchołkach,
osadzonym w $S^2$.
Wówczas brzeg każdej
ściany jest obrazem
cyklu prostego grafu $H$.
Jeśli nowy wierzchołek
leży we wnętrzu ściany
ograniczonej cyklem $C$,
można go połączyć
łukami parami rozłącznymi
poza nim z dowolnymi
różnymi wierzchołkami $C$,
bez przecinania
istniejącego rysunku.
\end{lemat}
\begin{proof}
Graf $H$ jest spójny i nie ma mostów: gdyby
krawędź $uv$ była mostem, po jej usunięciu powstałyby
dwie składowe zawierające odpowiednio $u$ i $v$.
Ponieważ graf ma co najmniej trzy wierzchołki,
jedna z tych składowych ma co najmniej dwa wierzchołki.
Usunięcie końca mostu leżącego w tej składowej
rozspójniałoby $H$, wbrew dwuspójności.
Z opisu regularnego otoczenia w dowodzie twierdzenia
\ref{thm:tg-euler-plane} wynika, że każdej ścianie
spójnego grafu na sferze odpowiada dokładnie jeden
obchód brzegowy. Z wniosku \ref{cor:tg-bridge-face}
i braku mostów wynika, że obie strony każdej krawędzi
należą do różnych ścian. Stąd żadna krawędź nie
występuje dwukrotnie w jednym obchodzie.
Przypuśćmy, że obchód powtarza wierzchołek $v$,
nie licząc jego powtórzenia przy zamknięciu obchodu. W dysku regularnego
otoczenia wokół $v$ dwa różne wystąpienia odpowiadają
dwóm różnym sektorom tej samej ściany $F$ pomiędzy
kolejnymi incydencjami krawędzi.

Wybierzmy w tych sektorach dwa krótkie promienie
wychodzące z $v$. Ich wnętrza leżą w $F$.
Ponieważ $F$ jest spójnym otwartym podzbiorem
sfery, jest łukowo spójny: wybierając punkt grafu
jako punkt w nieskończoności, utożsamiamy $F$ ze
spójnym otwartym podzbiorem płaszczyzny, w którym
punkty można łączyć łamanymi, a z łamanej usunąć
pętle, uzyskując łuk prosty.
Niech $\alpha:[0,1]\to F$ będzie takim łukiem
łączącym końce wybranych promieni. Niech $t_2$
będzie pierwszą chwilą spotkania drugiego domkniętego
promienia, a $t_1$ ostatnią chwilą spotkania pierwszego
promienia przed $t_2$. Czasy te istnieją dzięki
domkniętości promieni i ciągłości $\alpha$; są różne,
bo promienie przecinają się tylko w $v\notin F$.
Ograniczenie $\alpha$ do $[t_1,t_2]$ nie spotyka
żadnego z obu promieni poza końcami.
Wraz z odcinkami promieni do $v$ tworzy on
krzywą Jordana $\Gamma$ przecinającą $H$ tylko w $v$.
W obu cyklicznych przedziałach między wybranymi
sektorami znajduje się co najmniej jedna incydencja
krawędzi; ich łuki wychodzą z $v$ po przeciwnych
stronach $\Gamma$. Ponieważ $H$ jest prosty,
drugie końce tych krawędzi są różne od $v$.
Żadna droga w $H-v$ nie może połączyć tych końców,
gdyż musiałaby przeciąć $\Gamma$.
Sprzeczność z dwuspójnością dowodzi, że obchód
nie powtarza wierzchołków.

Jest on zatem cyklem prostym $C$.
Z twierdzenia Jordana--Schoenfliesa jego obraz
ogranicza dwa otwarte dyski na sferze.
Ściana $F$ leży w jednym z nich. Gdyby wewnątrz
tego dysku znajdowała się część grafu, spójność
$H$ dawałaby krawędź wychodzącą z pewnego
wierzchołka $C$ do tego dysku (lub cięciwę $C$).
Taka incydencja rozdzielałaby dwie kolejne
incydencje obchodu po stronie $F$, sprzecznie
z regułą następstwa w obchodzie ścianowym.
Zatem cały dysk jest ścianą $F$.
Za pomocą homeomorfizmu jego domknięcia
ze standardowym dyskiem możemy połączyć dowolny
punkt wnętrza z wybranymi różnymi punktami brzegu
odcinkami o rozłącznych wnętrzach. Dokładniej,
jeżeli obraz wybranego punktu nie jest środkiem
dysku jednostkowego, stosujemy automorfizm dysku
zespolonego przenoszący go do zera, przedłużający
się homeomorficznie na okrąg brzegowy.
Przenosząc je z powrotem, otrzymujemy żądane łuki.
\end{proof}


\begin{proof}[Dowód implikacji
\texorpdfstring{$(ii)\Rightarrow(i)$}{(ii) => (i)}
w twierdzeniu \ref{thm:tg-kuratowski}]
Załóżmy przeciwnie, że
istnieje nieplanarny
skończony graf prosty
niezawierający ani
subdywizji $K_5$, ani
subdywizji $K_{3,3}$.
Wśród wszystkich takich
grafów wybierzmy $G$,
dla którego
$|V(G)|+|E(G)|$
jest najmniejsze.
Wszystkie mniejsze
grafy bez wskazanych
subdywizji są więc
planarne. Ponieważ
każdy graf prosty
o co najwyżej czterech
wierzchołkach jest
podgrafem planarnego $K_4$,
mamy $|V(G)|\geq5$.

\emph{Krok 1: graf $G$
jest dwuspójny.}
Gdyby był niespójny,
jego składowe byłyby
mniejszymi grafami,
a więc planarnymi;
można by je narysować
w rozłącznych dyskach,
otrzymując osadzenie $G$.
Gdyby miał wierzchołek
rozcinający, moglibyśmy
podzielić go na właściwe
podgrafy $G_1,G_2$
mające tylko ten jeden
wspólny wierzchołek
i pozbawione krawędzi
między pozostałymi
częściami. Oba są
mniejsze i nie
zawierają zabronionych
subdywizji, więc
są planarne.
Lemat
\ref{lem:tg-kuratowski-gluing}(i)
uczyniłby planarnym
także $G$,
co jest sprzecznością.

\emph{Krok 2: graf $G$
jest trójspójny.}
Przypuśćmy, że
istnieją różne
wierzchołki $x,y$,
których usunięcie
rozspójnia $G$.
Wybierzmy niepusty
podział składowych
$G-\{x,y\}$ na dwie
rodziny. Niech
$G_1,G_2$ będą
podgrafami złożonymi
odpowiednio z jednej
i drugiej rodziny
oraz wierzchołków
$x,y$, a istniejącą
krawędź $xy$ (o ile
jest obecna)
przypiszmy do $G_1$.
Wspólne wierzchołki
tych grafów to
dokładnie $x,y$.

Ponieważ $G$ jest
dwuspójny, każda
składowa $G-\{x,y\}$
ma sąsiada zarówno
w $x$, jak i w $y$;
inaczej jeden z
$x,y$ rozcinałby $G$.
Istnieje więc
w każdej z dwóch
części ścieżka
$x$--$y$ o wszystkich
wewnętrznych
wierzchołkach
w tej części.
Oznaczmy takie
ścieżki przez $Q_1,Q_2$.
Niech $H_i=G_i+xy$.
Każdy $H_i$ ma
mniej wierzchołków
niż $G$, zaś
co najwyżej jedną
krawędź, której
nie było w $G_i$.
Co najmniej dwie
krawędzie $G$
leżą poza $G_i$
(ścieżka $Q_{3-i}$
ma długość co
najmniej dwa).
Zatem
\[
|V(H_i)|+|E(H_i)|
<|V(G)|+|E(G)|.
\]

Gdyby w $H_i$
istniała zabroniona
subdywizja, w której
wykorzystano nowo
dodane $xy$, zastępujemy
w niej ten jeden
łuk ścieżką
$Q_{3-i}$. Jej
wnętrze jest rozłączne
z $G_i$, więc
otrzymujemy tę samą
subdywizję w $G$.
Jeżeli subdywizja
nie używa $xy$ albo
$xy$ już należy do
$G$, jest już
podgrafem $G$.
Sprzeczność
pokazuje, że
oba $H_i$ są
wolne od zabronionych
subdywizji, a przez
minimalność $G$
są planarne.
Lemat
\ref{lem:tg-kuratowski-gluing}(ii)
daje wtedy
planarność $G$,
kolejną sprzeczność.
To dowodzi trójspójności.

\emph{Krok 3: kurczymy
odpowiednią krawędź.}
Ponieważ $G\not\cong K_4$,
z lematu
\ref{lem:tg-kuratowski-contractible-edge}
istnieje krawędź $xy$
taka, że prosty graf
$G'=G/xy$ jest
trójspójny.
Niech $v$ będzie
wierzchołkiem powstałym
z kontrakcji.
Z wniosku
\ref{cor:tg-kuratowski-minor-heredity}
graf $G'$ nie zawiera
żadnej zabronionej
subdywizji.
Kontrakcja zmniejsza liczbę wierzchołków o jeden
i nie zwiększa liczby krawędzi: usuwamy $xy$,
a ewentualne krawędzie równoległe zastępujemy
pojedynczymi. Zatem
$|V(G')|+|E(G')|<|V(G)|+|E(G)|$.
Z minimalności kontrprzykładu wynika planarność $G'$.

Ustalmy osadzenie
$G'$ na sferze.
Graf $G'-v$ jest
dwuspójny, ponieważ
$G'$ jest trójspójny.
Niech $F$ będzie
ścianą osadzenia
$G'-v$ zawierającą
wcześniejszy punkt $v$.
Na mocy lematu
\ref{lem:tg-kuratowski-face-cycle}
jej brzeg jest
cyklem prostym $C$.
Po usunięciu $v$ i jego incydentnych krawędzi
wnętrza tych krawędzi wraz z punktem $v$ tworzą
spójny podzbiór dopełnienia osadzenia $G'-v$.
Należą więc do jednej ściany $F$. Ich drugie końce
leżą w $G'-v$ i w domknięciu $F$, a zatem na $C$.
Oznaczmy
\[
X=N_G(x)\setminus\{y\},
\qquad
Y=N_G(y)\setminus\{x\}.
\]
Są one zawarte
w $V(C)$.
Każdy wierzchołek grafu $G$ ma stopień co najmniej
trzy. W przeciwnym razie usunięcie wszystkich jego
sąsiadów (co najwyżej dwóch) izolowałoby go od
pozostałych wierzchołków, których istnienie wynika
z $|V(G)|\geq5$. Ponieważ $xy\in E(G)$, otrzymujemy
$|X|\geq2$ i $|Y|\geq2$.

Niech $x_1,\ldots,x_k$
będą kolejnymi
wierzchołkami $X$
w porządku cyklicznym
na $C$, $k\geq2$.
Przez $C_i$
oznaczmy łuk cyklu
$C$ od $x_i$
do $x_{i+1}$
niezawierający
w swoim wnętrzu
żadnego wierzchołka $X$
(indeksy modulo $k$).
Przy $k=2$ są to
dwa różne łuki
o tych samych końcach.

\emph{Krok 4:
możliwość wstawienia $y$.}
Gdyby istniało $i$
takie, że
$Y\subseteq V(C_i)$,
można by w osadzeniu
$G'$ zostawić
wierzchołek $v$
jako $x$, usunąć
ewentualne krawędzie
do sąsiadów z
$Y\setminus X$
i zachować
jego łuki do
wszystkich $x_j$.
Zachowane łuki od $x$ do $x_j$ leżą wewnątrz
domkniętego dysku $\overline F$ poza końcami
na $C$; ich jedynym wspólnym punktem jest $x$.
Łuki te dzielą $F$ na otwarte dyski: każdy z nich
jest ograniczony przez łuki $xx_i$, $xx_{i+1}$
i ścieżkę $C_i$. Istotnie, suma tych trzech łuków
jest krzywą Jordana, a po stronie przylegającej
do wnętrza $C_i$ nie ma innego łuku wychodzącego
z $x$: musiałby on albo kończyć się we wnętrzu
$C_i$ w punkcie $X$, albo przeciąć krzywą.
Obie możliwości są wykluczone. Dotyczy to także
przypadku $k=2$, kiedy otrzymujemy dwa dyski.
W dysku odpowiadającym $C_i$ umieszczamy $y$.
Punkty $x$ i wszystkie wierzchołki $Y$ są różnymi
punktami jego brzegu, ponieważ $Y\subseteq V(C_i)$.
Przenosząc domknięty dysk na dysk standardowy
i rysując odcinki od środka do tych punktów,
jak w ostatniej części dowodu lematu
\ref{lem:tg-kuratowski-face-cycle}, uzyskujemy łuki
rozłączne poza wspólnym końcem $y$ i nieprzecinające
dotychczasowego grafu poza ich drugimi końcami.
Odtwarza to
wszystkie krawędzie
grafu $G$,
bez przecięć.
Sprzeczność.
Dlatego
\[
Y\not\subseteq V(C_i)
\quad\text{dla każdego }i.
\tag{*}
\]

\emph{Krok 5:
trzech wspólnych sąsiadów.}
Przypuśćmy, że
$|X\cap Y|\geq3$.
Niech $a,b,c$
będą trzema różnymi
wspólnymi sąsiadami
$x,y$ leżącymi na $C$.
Po ewentualnej zmianie nazw przyjmujemy, że $a,b,c$
występują w tej kolejności na zorientowanym cyklu $C$.
Dzielimy zbiór wierzchołków cyklu na trzy przedziały
cykliczne: od $a$ do poprzednika $b$, od $b$ do
poprzednika $c$ i od $c$ do poprzednika $a$.
Każdy z nich jest niepusty (może być singletonem)
i wyznacza spójną ścieżkę na $C$. Ścieżki są
parami wierzchołkowo rozłączne, zawierają odpowiednio
$a,b,c$ i są połączone trzema pominiętymi
krawędziami cyklu.
Wraz z singletonami
$\{x\}$ i $\{y\}$
są to pięć zbiorów
rozgałęzienia:
trzy ścieżki tworzą
trójkąt po kontrakcji,
a każdy z $x,y$
jest połączony
z każdą z nich;
dodatkowo $xy\in E(G)$.
Otrzymujemy
model minora $K_5$ w $G$. Na mocy lematu
\ref{lem:tg-kuratowski-model-minor} graf $K_5$ jest
minorem $G$, więc z lematu
\ref{lem:tg-kuratowski-minor-to-subdivision}
graf $G$ zawiera subdywizję $K_5$ lub $K_{3,3}$.
To przeczy
założeniu o $G$.
Wobec tego
\[
|X\cap Y|\leq2.
\tag{**}
\]

\emph{Krok 6:
naprzemienne końce
i subdywizja $K_{3,3}$.}
Wskażemy cztery
różne wierzchołki
cyklu $C$ w porządku
cyklicznym
$y_1,x',y_2,x''$
takie, że
$y_1,y_2\in Y$,
$x',x''\in X$.

Jeśli $Y\setminus X$
jest niepuste,
wybierzmy $y_1$
w tym zbiorze.
Leży on we wnętrzu
pewnego $C_i$.
Z warunku (*) istnieje
$y_2\in Y\setminus V(C_i)$.
Wtedy cztery
różne punkty
$y_1,x_{i+1},y_2,x_i$
leżą na $C$
w odpowiednim
cyklicznym porządku;
wybieramy
$x'=x_{i+1}$,
$x''=x_i$.

Jeśli zaś $Y\subseteq X$,
to na mocy (**) i
$|Y|\geq2$
mamy $Y=\{y_1,y_2\}$.
Warunek (*) oznacza,
że te dwa wierzchołki
nie są kolejnymi
elementami $X$
w porządku cyklicznym:
w przeciwnym razie
leżałyby na jednym
łuku $C_i$.
Na każdym z dwóch
łuków cyklu
między $y_1$ i $y_2$
istnieje więc
wierzchołek $X$
różny od końców.
Wybierzmy po jednym
z nich jako $x',x''$.
Uzyskujemy ten sam
naprzemienny porządek.

Cztery łuki cyklu $C$
między kolejnymi
wierzchołkami
$y_1,x',y_2,x''$,
wraz z krawędziami
$xy$, $xx'$, $xx''$,
$yy_1$, $yy_2$,
tworzą podgraf
izomorficzny z
subdywizją $K_{3,3}$.
Jego dwie części
rozgałęzienia to
\[
\{x,y_1,y_2\},
\qquad
\{y,x',x''\}.
\]
Istotnie, $x$ łączy się
z trzema wierzchołkami
drugiej części
przez $xy,xx',xx''$;
każdy z $y_1,y_2$
łączy się z $y$
odpowiednią krawędzią
i z $x',x''$
przez dwa rozłączne
łuki cyklu $C$.
Wnętrza czterech
łuków po $C$
i wymienionych
pięciu krawędzi
są parami rozłączne:
łuki na $C$
stykają się tylko
w końcach, a $x,y$
nie należą do $C$.
Otrzymujemy
zabronioną subdywizję.

Wszystkie możliwości
prowadzą do sprzeczności.
Zatem nie istnieje
nieplanarny graf prosty
bez subdywizji $K_5$
i $K_{3,3}$.
Jest to implikacja
$(ii)\Rightarrow(i)$,
która wraz z
udowodnionym wcześniej
kierunkiem
$(i)\Rightarrow(ii)$
kończy dowód twierdzenia
Kuratowskiego.
\end{proof}

\begin{uwaga}[Zależności dowodowe]
W powyższym dowodzie
nie przyjęliśmy
twierdzenia Kuratowskiego
ani Wagnera jako lematu.
Wykorzystaliśmy twierdzenie Jordana--Schoenfliesa
oraz lokalne prostowanie skończonych grafów
w powierzchni i opis regularnych otoczeń, przyjęte
w dowodzie twierdzenia \ref{thm:tg-euler-plane}.
Korzystaliśmy także z elementarnych własności
spójnych otwartych podzbiorów płaszczyzny i z
wniosku \ref{cor:tg-bridge-face} o mostach.
Żaden z tych wyników nie wymaga kryteriów
planarności Kuratowskiego ani Wagnera.
Metoda redukcji przez
kontrakcję zachowującą
trójspójność jest
klasycznie związana
z dowodem Thomassena
(zob. także R.~Diestel,
\emph{Graph Theory},
rozdział o planarności).
Niezależność twierdzenia
topologicznego od
kryteriów grafowych
zapobiega błędnemu kołu.
\end{uwaga}

\begin{twierdzenie}[Wagner -- sformułowanie przez minory]
\label{thm:tg-wagner}
Skończony graf prosty $G$ jest planarny wtedy i tylko wtedy,
gdy nie zawiera $K_5$ ani $K_{3,3}$ jako minora.
\end{twierdzenie}

\begin{proof}
Jeśli $G$ jest planarny, to nie może mieć
nieplanarnego minora: dowód zamknięcia planarności
na minory podaliśmy powyżej,
a $K_5$ i $K_{3,3}$ nie są planarne.

Odwrotnie, załóżmy, że $G$ nie ma żadnego
z tych dwóch minorów. Gdyby zawierał
subdywizję $K_5$ lub $K_{3,3}$,
kontrakcja jej dróg odpowiadających
krawędziom grafu wyjściowego
dawałaby taki minor w $G$.
Zatem $G$ nie zawiera żadnej
z tych dwóch zabronionych subdywizji.
Twierdzenie \ref{thm:tg-kuratowski},
którego oba kierunki już udowodniliśmy,
gwarantuje zatem planarność $G$.
\end{proof}

\begin{uwaga}
Relacje bycia minorem i bycia podgrafem topologicznym
nie są równoważne dla dowolnej pary grafów.
Lemat \ref{lem:tg-kuratowski-minor-to-subdivision}
pozwala jednak przejść między odpowiednimi
\emph{parami zabronionych przeszkód}:
$K_5$ może przy tym zostać zastąpiony
przez subdywizję $K_{3,3}$.
Z tego powodu twierdzenie Wagnera
jest konsekwencją pełnego twierdzenia Kuratowskiego,
a nie dodatkową przesłanką w jego dowodzie.
\end{uwaga}

\begin{przyklad}[Dlaczego sama nierówność Eulera nie wystarcza]
Dla grafu $K_{3,3}$ mamy $n=6$ i $m=9$.
Ograniczenie $m\le3n-6$ daje tylko $9\le12$ i nie wykrywa
nieplanarności. Dopiero wykorzystanie dwudzielności,
a więc braku cykli nieparzystych, daje sprzeczność
$9\le2\cdot6-4=8$. Dla innych grafów mogą być potrzebne
jeszcze silniejsze narzędzia, takie jak twierdzenie Kuratowskiego.
\end{przyklad}

\section{Grafy zewnętrznie planarne}
\label{sec:tg-outerplanar}

\begin{definicja}[Zewnętrzna planarność]
Graf skończony nazywamy \emph{zewnętrznie planarnym}, jeśli
ma zanurzenie w płaszczyźnie, w którym każdy wierzchołek
leży na brzegu ściany nieograniczonej.
\end{definicja}

\begin{twierdzenie}[Kryterium zewnętrznej planarności przez dodanie wierzchołka]
\label{thm:tg-outerplanar-cone}
Niech $G$ będzie skończonym grafem prostym.
Przez $G^+$ oznaczmy graf powstały przez dodanie
nowego wierzchołka $p$ i krawędzi $pv$ dla każdego
$v\in V(G)$. Wówczas $G$ jest zewnętrznie planarny
wtedy i tylko wtedy, gdy $G^+$ jest planarny.
\end{twierdzenie}

\begin{proof}
Jeżeli $G$ nie ma wierzchołków, teza wynika
bezpośrednio z definicji: $G$ jest zewnętrznie
planarny, a $G^+=K_1$ jest planarny.
Najpierw rozważmy zatem spójny graf zewnętrznie
planarny o co najmniej dwóch wierzchołkach. Wybieramy regularne otoczenie
w sferze jak w dowodzie twierdzenia \ref{thm:tg-euler-plane}.
Zewnętrznej ścianie odpowiada jeden okrąg brzegowy
tego otoczenia; po przeciwnej stronie otoczenia
okrąg ogranicza dysk. Każdy wierzchołek grafu ma
co najmniej jeden sektor przyległy do zewnętrznej
ściany. Wybieramy po jednym takim sektorze
i prowadzimy w nim krótki łuk od wierzchołka
do odpowiedniego punktu okręgu brzegowego.
Łuki te są rozłączne, gdyż leżą w różnych dyskach
wierzchołkowych. Umieszczamy $p$ we wnętrzu
dysku po drugiej stronie okręgu i łączymy go
z wybranymi punktami brzegu łukami rozłącznymi
poza $p$, korzystając z homeomorfizmu z dyskiem
standardowym. Po przedłużeniu przez wybrane
sektory otrzymujemy wszystkie krawędzie $pv$.
Dla składowej jednoelementowej jej rozszerzenie
jest po prostu grafem $K_2$, więc jest planarne.
Jeżeli $G$ jest niespójny, każda jego składowa
dziedziczy zewnętrzną planarność: po usunięciu
pozostałych składowych dotychczasowa ściana
nieograniczona zawiera się w nowej ścianie
nieograniczonej, a zachowane wierzchołki nadal
należą do jej brzegu. Stosujemy powyższą konstrukcję
do każdej składowej, używając tej samej nazwy $p$
dla nowego wierzchołka. Sklejamy otrzymane grafy
kolejno w $p$ na mocy lematu
\ref{lem:tg-kuratowski-gluing}; ich suma jest
dokładnie grafem $G^+$.

Odwrotnie, osadźmy $G^+$ w sferze i usuńmy $p$
wraz z incydentnymi krawędziami. Punkt $p$ oraz
wnętrza wszystkich usuniętych krawędzi leżą
w jednej składowej dopełnienia obrazu $G$,
ponieważ tworzą zbiór spójny. Każdy wierzchołek
$G$ leży w domknięciu tej składowej jako drugi
koniec jednej z usuniętych krawędzi, ale nie
należy do niej. Wybieramy punkt $q$ w tej ścianie i utożsamiamy
$S^2\setminus\{q\}$ z płaszczyzną. Usunięcie punktu
z otwartego spójnego podzbioru powierzchni nie
rozspaja go, jak wykazano w rozdziale
\ref{ch:tg-planarne}. Rozpatrywana ściana staje się
więc ścianą nieograniczoną, a wszystkie zachowane
wierzchołki pozostają na jej brzegu.
\end{proof}

\begin{twierdzenie}[Zakazane minory grafów zewnętrznie planarnych]
\label{thm:tg-outerplanar-minors}
Skończony graf prosty $G$ jest zewnętrznie planarny
wtedy i tylko wtedy, gdy nie ma minora izomorficznego
z $K_4$ ani z $K_{2,3}$.
\end{twierdzenie}

\begin{proof}
Jeżeli $H$ jest minorem $G$, to model $H$ w $G$
wraz z dodatkowym singletonem $\{p\}$ daje model
$H^+$ w $G^+$. Ponieważ $K_4^+=K_5$, a $K_{2,3}^+$
zawiera $K_{3,3}$ jako podgraf (dodajemy $p$ do
dwuelementowej części), istnienie któregokolwiek
z tych minorów przeczy planarności $G^+$ na mocy
twierdzenia Wagnera \ref{thm:tg-wagner}.

Odwrotnie, jeżeli $G$ nie jest zewnętrznie planarny,
to $G^+$ nie jest planarny. Ma więc model $K_5$
lub $K_{3,3}$. Co najwyżej jeden zbiór rozgałęzienia
zawiera $p$. W przypadku $K_5$ odrzucamy ten
zbiór, a jeśli żaden nie zawiera $p$, odrzucamy
dowolny. Pozostałe cztery zbiory tworzą model $K_4$ w $G$:
żaden z nich nie zawiera $p$, więc ich podgrafy
indukowane oraz świadkowie krawędzi między nimi
są takie same w $G^+$ i w $G$. W przypadku $K_{3,3}$ odrzucamy
zbiór zawierający $p$ lub dowolny jeden zbiór,
jeśli $p$ nie występuje w modelu. Pozostałe zbiory tworzą model $K_{2,3}$ w $G$
(lub $K_{3,2}$, izomorficznego z $K_{2,3}$).
Ponownie spójność i wszystkie wymagane krawędzie
pozostają zachowane po usunięciu $p$.
To dowodzi równoważności.
\end{proof}

\begin{twierdzenie}[Liczba krawędzi grafu zewnętrznie planarnego]
Niech $G$ będzie prostym, spójnym grafem zewnętrznie planarnym,
o $n\ge2$ wierzchołkach i $m$ krawędziach. Wtedy
$m\le2n-3$. Dla $n\ge3$ równość zachodzi wtedy i tylko
wtedy, gdy w ustalonym osadzeniu zewnętrznie planarnym
brzeg ściany zewnętrznej jest cyklem prostym
zawierającym wszystkie wierzchołki, a wszystkie
ściany ograniczone mają brzegi trójkątne.
\end{twierdzenie}

\begin{proof}
Ustalmy zanurzenie z wszystkimi wierzchołkami na brzegu
ściany zewnętrznej. Jeśli $n=2$, mamy co najwyżej jedną
krawędź i wynik jest oczywisty. Przyjmijmy $n\ge3$.
Najpierw uzupełnimy graf do grafu dwuspójnego bez zmiany
zbioru wierzchołków. Jeśli istnieje wierzchołek rozcinający
$v$, obchód brzegu zewnętrznego odwiedza $v$ co najmniej
dwukrotnie. Każda składowa $G-v$ zawiera wierzchołki występujące
w tym obchodzie. Ponieważ obchód jest cykliczny,
istnieje wystąpienie $v$ pomiędzy krawędziami $av$ i $vb$,
dla których $a,b$ należą do różnych składowych $G-v$:
inaczej przy przejściu przez $v$ nigdy nie zmieniałaby
się składowa, a obchód nie mógłby odwiedzać ich wszystkich.
W szczególności $a\ne b$ i $ab\notin E(G)$.
W sektorze ściany zewnętrznej odpowiadającym temu
wystąpieniu prowadzimy łuk $ab$ blisko drogi $a-v-b$.
Odcina on od ściany zewnętrznej mały obszar ograniczony
przez $av$, $vb$ i nowy łuk, niezawierający innych
fragmentów grafu. Wierzchołki $a,b$ pozostają na brzegu
ściany zewnętrznej przy końcach łuku, a $v$ pozostaje
na nim dzięki innemu wystąpieniu w obchodzie.
Pozostałe wierzchołki nie zmieniają swego położenia
względem ściany zewnętrznej.
Po dodaniu $ab$ liczba składowych grafu po usunięciu $v$
maleje o jeden. Dla każdego innego wierzchołka $u$
liczba składowych po usunięciu $u$ nie wzrasta, gdyż
dodanie krawędzi nie może rozspójnić grafu.
Nieujemna całkowita suma
$\sum_{u\in V(G)}(c(G-u)-1)$, gdzie $c$ oznacza
liczbę składowych, maleje więc ściśle. Proces kończy
się przy braku wierzchołków rozcinających. Otrzymany
graf jest dwuspójny ($n\ge3$), nadal zewnętrznie
planarny i ma ten sam zbiór wierzchołków.
Na mocy lematu \ref{lem:tg-kuratowski-face-cycle}
w grafie dwuspójnym obchód ściany zewnętrznej jest
cyklem prostym. Ponieważ każdy wierzchołek leży na brzegu
zewnętrznym, jest to cykl długości $n$.

Każda pozostała ściana jest ograniczona cyklem prostym.
Jeżeli ma co najmniej cztery boki, wybieramy cztery
kolejne wierzchołki $a,b,c,d$ jej brzegu.
Przynajmniej jednej z krawędzi $ac,bd$ jeszcze nie ma:
gdyby obie istniały, ich łuki leżałyby po drugiej
stronie cyklu brzegowego ściany i miały naprzemienne
końce na brzegu dysku, więc z twierdzenia Jordana
musiałyby się przeciąć. Brakującą przekątną rysujemy
wewnątrz otwartego dysku będącego ścianą. Nie powstaje
przecięcie ani zmiana zewnętrznego obwodu. Operację
powtarzamy, aż każda wewnętrzna ściana będzie trójkątem.
Ponieważ za każdym razem zwiększamy liczbę krawędzi,
a graf pozostaje prosty na $n$ wierzchołkach, proces
kończy się po skończonej liczbie kroków.

Oznaczmy przez $m'$ liczbę krawędzi grafu uzupełnionego,
a przez $f$ liczbę jego ścian na sferze.
Suma długości ich obwodów wynosi $2m'$, przy czym
ściana zewnętrzna ma obwód długości $n$,
a każda z $f-1$ ścian wewnętrznych długości trzy.
Stąd $2m'=n+3(f-1)$.
Z równania Eulera $f=2-n+m'$ dostajemy
$2m'=n+3(1-n+m')$, czyli $m'=2n-3$.
Ponieważ $G$ jest podgrafem grafu uzupełnionego,
mamy $m\leq m'=2n-3$.
Jeżeli $n\ge3$ oraz $m=2n-3$, żadna z operacji
uzupełniania nie mogła być wykonana: każda dodaje
krawędź, a końcowa liczba krawędzi wynosi dokładnie
$2n-3$. Zatem już pierwotne osadzenie ma opisany
cykl zewnętrzny i trójkątne ściany ograniczone.
Odwrotnie, jeżeli osadzenie spełnia te warunki,
to bez uzupełniania mamy $2m=n+3(f-1)$, skąd
$m=2n-3$.
\end{proof}

\begin{uwaga}[Bezpośrednie oszacowanie]
Nierówność można również otrzymać bez uzupełniania
grafu. Dla spójnego grafu o $n\geq3$ oznaczmy przez
$d$ długość obchodu zewnętrznej ściany. Ponieważ każdy
wierzchołek należy do jej brzegu, $d\geq n$.
Pozostałe ściany mają obchody długości co najmniej trzy
na mocy lematu o stopniach ścianowych z rozdziału
\ref{ch:tg-planarne}. Zatem $2m\geq n+3(f-1)$.
Ze wzoru Eulera dostajemy $2m\geq n+3(1-n+m)$,
czyli $m\leq2n-3$. Argument obowiązuje także wtedy,
gdy w zewnętrznym obchodzie powtarzają się wierzchołki.
\end{uwaga}

\begin{uwaga}
Dowód twierdzenia wykazuje również, że każde spójne
zewnętrznie planarne osadzenie grafu prostego o co
najmniej trzech wierzchołkach można uzupełnić do
triangulacji dysku ograniczonego cyklem zewnętrznym,
bez zmiany zbioru wierzchołków. Otrzymany graf jest
maksymalny w klasie grafów zewnętrznie planarnych
na tym zbiorze wierzchołków: dodanie kolejnej krawędzi
przeczyłoby wykazanemu ograniczeniu $m\le2n-3$. Samo liczenie
$2m=n+3(f-1)$ dotyczy dopiero grafu \emph{uzupełnionego},
a nie dowolnego grafu zewnętrznie planarnego.
\end{uwaga}

\section{Kolorowanie wierzchołków: twierdzenie o pięciu kolorach}
\label{sec:tg-five-color}

\begin{definicja}[Poprawne kolorowanie]
Niech $G=(V,E)$ będzie skończonym grafem prostym,
a $q\ge1$ liczbą całkowitą.
\emph{Poprawnym $q$-kolorowaniem wierzchołków} nazywamy
funkcję $c:V\to\{1,\ldots,q\}$, dla której
$c(u)\ne c(v)$ przy każdej krawędzi $uv\in E$.
\emph{Liczbą chromatyczną} $\chi(G)$ nazywamy najmniejsze
takie $q$, jeśli $V\ne\varnothing$ (istnieje ono,
gdyż wystarczy przypisać każdemu wierzchołkowi osobny
kolor). Dla grafu o co najmniej jednym wierzchołku
bez krawędzi mamy $\chi(G)=1$, ponieważ wszystkie
wierzchołki mogą otrzymać ten sam kolor.
Jedynie dla grafu zerowego $G_0$ przyjmujemy $\chi(G_0)=0$.
Przykładowo $\chi(\overline{K_4})=1$, natomiast
$\chi(K_4)=4$; ponadto $\chi(G_0)=0$.
Nie należy mylić liczby chromatycznej z liczbą wszystkich
poprawnych kolorowań: dla $\overline{K_4}$ jest ich $q^4$
przy $q$ oznaczonych kolorach.
\end{definicja}

\begin{lemat}[Wierzchołek stopnia co najwyżej dwa]
\label{lem:tg-outerplanar-degree-two}
Niech $G$ będzie skończonym prostym grafem zewnętrznie
planarnym o co najmniej jednym wierzchołku. Wtedy
$G$ ma wierzchołek stopnia co najwyżej dwa.
\end{lemat}

\begin{proof}
Wystarczy rozważyć jedną składową, ponieważ każda
składowa grafu zewnętrznie planarnego sama ma
zanurzenie zewnętrznie planarne. Dla składowej
o jednym lub dwóch wierzchołkach wynik jest oczywisty.
Niech zatem składowa ma $n\ge3$ wierzchołków.
Jak w dowodzie oszacowania $m\le2n-3$,
uzupełniamy ją do triangulacji wielokąta o $n$
wierzchołkach: krawędzie zewnętrzne tworzą prosty
cykl, a wszystkie ściany wewnętrzne są trójkątami.

Rozważmy graf pomocniczy, którego wierzchołki
odpowiadają trójkątom triangulacji, a krawędzie
przekątnym wspólnym dwóm trójkątom.
Ten graf jest spójny: drogę między punktami
we wnętrzach dwóch trójkątów można poprowadzić
wewnątrz wielokąta, omijając wierzchołki triangulacji
i przecinając przekątne poprzecznie.
Graf uzupełniony ma $2n-3$ krawędzi, w tym $n$
krawędzi obwodu i $n-3$ przekątnych. Ze wzoru Eulera
ma $n-2$ ścian wewnętrznych. Każda przekątna graniczy
z dwoma różnymi trójkątami, a różne trójkąty mają
co najwyżej jeden wspólny bok. Graf pomocniczy ma
zatem $n-2$ wierzchołki i $n-3$ krawędzie, więc jest drzewem. Dla $n=3$ jedyny trójkąt
daje wprost wierzchołek stopnia dwa.
Dla $n>3$ wybieramy liść drzewa pomocniczego.
Odpowiada mu trójkąt mający tylko jeden bok
wewnętrzny; jego pozostałe dwa boki leżą na
zewnętrznym obwodzie. Ich wspólny wierzchołek
nie jest końcem żadnej dodatkowej przekątnej,
gdyż przy tym wierzchołku sektor wewnątrz wielokąta
pomiędzy dwoma bokami obwodu jest w całości sektorem
wybranego trójkąta. Każda dodatkowa przekątna musiałaby
wychodzić właśnie przez ten sektor, a więc przecinać
wnętrze ściany.
Ma zatem stopień dwa w grafie uzupełnionym,
a w pierwotnym grafie stopień nie większy niż dwa.
\end{proof}

\begin{twierdzenie}[Trzy kolory dla grafów zewnętrznie planarnych]
\label{thm:tg-outerplanar-three-color}
Każdy skończony prosty graf zewnętrznie planarny
można poprawnie pokolorować co najwyżej
trzema kolorami.
\end{twierdzenie}

\begin{proof}
Prowadzimy indukcję względem liczby wierzchołków.
Dla grafu pustego istnieje jedyne kolorowanie pustą
funkcją. Dla grafu niepustego lemat
\ref{lem:tg-outerplanar-degree-two} zapewnia
wierzchołek $v$ o stopniu nie większym niż dwa.
Graf $G-v$ jest zewnętrznie planarny: na mocy
twierdzenia \ref{thm:tg-outerplanar-minors}
brak zakazanych minorów przechodzi na podgraf,
gdyż minor podgrafu jest również minorem $G$.
Zatem z założenia indukcyjnego $G-v$ ma poprawne
trzykolorowanie. Co najwyżej dwa kolory
występują na sąsiadach $v$; wybieramy
jeden z pozostałych kolorów dla $v$.
Otrzymujemy poprawne trzykolorowanie całego $G$.
\end{proof}



\begin{lemat}[Wierzchołek małego stopnia]
Każdy skończony prosty graf planarny o co najmniej jednym
wierzchołku ma wierzchołek stopnia co najwyżej pięć.
\end{lemat}

\begin{proof}
Dla $n\le2$ teza jest oczywista. Jeśli $n\ge3$, możemy
najpierw dołożyć krawędzie łączące składowe w planie,
uzyskując graf spójny i prosty. Twierdzenie
\ref{thm:tg-planarne-bound} daje $2m\le6n-12<6n$.
Z lematu o uściskach dłoni średni stopień wynosi $2m/n<6$.
Zatem co najmniej jeden stopień jest nie większy niż pięć.
\end{proof}

\begin{twierdzenie}[Twierdzenie o pięciu kolorach]
\label{thm:tg-5color}
Każdy skończony graf prosty planarny można poprawnie
pokolorować co najwyżej pięcioma kolorami.
\end{twierdzenie}

\begin{proof}
Stosujemy indukcję po liczbie wierzchołków. Dla
$n\le5$ twierdzenie jest oczywiste. Niech $G$
będzie planarny o $n>5$ wierzchołkach. Z lematu wybieramy
wierzchołek $v$ stopnia co najwyżej pięć. Z założenia
indukcyjnego graf $G-v$ ma poprawne pięciokolorowanie.
Jeżeli sąsiedzi $v$ używają co najwyżej czterech kolorów,
możemy nadać $v$ jeden z niewykorzystanych kolorów.
Rozważmy zatem przypadek, w którym $v$ ma pięciu
różnych sąsiadów i każdy z nich ma inny kolor.
Ustalamy osadzenie $G$ w płaszczyźnie. Oznaczamy
sąsiadów przez $v_1,\ldots,v_5$ zgodnie z porządkiem
cyklicznym krawędzi przy $v$ i permutujemy nazwy
kolorów tak, aby $v_i$ miał kolor $i$.

Rozważmy podgraf $H_{13}$ grafu $G-v$ indukowany przez
wierzchołki o kolorach $1$ i $3$. Jeżeli $v_1$ i $v_3$
należą do różnych składowych $H_{13}$, zamieniamy
kolory $1$ i $3$ na składowej zawierającej $v_1$.
Na krawędziach o obu końcach w tej składowej
zamiana zachowuje różność kolorów. Jeżeli tylko jeden
koniec krawędzi należy do składowej, drugi koniec
nie ma koloru $1$ ani $3$: w przeciwnym razie należałby
do tej samej składowej $H_{13}$. Krawędzie o obu końcach
poza składową nie zmieniają kolorów końców.
Otrzymujemy więc nadal poprawne kolorowanie.
Jedyny sąsiad $v$ mający wcześniej kolor $1$, czyli
$v_1$, ma teraz kolor $3$, a $v_3$ pozostaje poza
składową. Kolor $1$ znika z sąsiedztwa $v$;
kolorujemy wtedy $v$ kolorem $1$.

W przeciwnym przypadku wybieramy prostą ścieżkę
w $H_{13}$ łączącą $v_1$ z $v_3$. Istnieje ona, gdyż
w skończonym grafie spójnym można usunąć powtórzenia
wierzchołków z dowolnego łączącego je spaceru.
Ścieżka nie zawiera $v$; po dodaniu krawędzi
$vv_1$ i $vv_3$ otrzymujemy cykl prosty, którego obraz
w ustalonym osadzeniu jest krzywą Jordana.
W ustalonym zanurzeniu planarnym sąsiedzi $v_2$ i $v_4$
leżą po przeciwnych stronach tej krzywej, ponieważ
ich krawędzie do $v$ wychodzą pomiędzy dwiema wskazanymi
krawędziami w różnych sektorach lokalnego dopełnienia
krzywej. Wnętrza tych krawędzi i ich drugie końce
nie spotykają krzywej, więc każdy z tych łuków
pozostaje w jednej z dwóch składowych jej dopełnienia.
Zatem nie może istnieć ścieżka złożona z wierzchołków
kolorów $2$ i $4$, łącząca $v_2$ z $v_4$:
jej obraz musiałby przeciąć poprzednią krzywą.
Jednak ścieżka w grafie $G-v$ używająca wyłącznie
kolorów $2,4$ nie ma wspólnego wierzchołka z cyklem
zbudowanym z $v$ i wierzchołków kolorów $1,3$.
W osadzeniu grafu obrazy tych podgrafów są rozłączne.
Oznaczmy przez $H_{24}$ podgraf grafu $G-v$ indukowany
przez wierzchołki o kolorach $2$ i $4$. Wierzchołki
$v_2,v_4$ należą do różnych jego składowych.
Zamieniamy kolory $2$ i $4$ w składowej $H_{24}$
zawierającej $v_2$. Ten sam argument dotyczący
krawędzi na granicy składowej zapewnia poprawność
nowego kolorowania. Wierzchołek $v_4$ nie zmienia
koloru, więc kolor $2$ znika z sąsiedztwa $v$
i można go przypisać wierzchołkowi $v$.
W obu przypadkach przedłużamy kolorowanie, co kończy indukcję.
\end{proof}

\begin{center}
\begin{tikzpicture}[scale=.9,line cap=round]
  \node[tgVertex] (v) at (0,0) {$v$};
  \foreach \i/\ang/\col in {1/90/tgBlue,2/18/tgTeal,3/-54/tgCoral,4/-126/tgGold,5/162/tgPurple} {
    \node[tgVertex,draw=\col,fill=\col!15] (n\i) at (\ang:2.0) {$\i$};
    \draw[tgEdge] (v)--(n\i);
  }
  \node[tgVertex,draw=tgCoral,fill=tgCoral!15] (a) at (60:3.8) {$3$};
  \node[tgVertex,draw=tgBlue,fill=tgBlue!15] (b) at (-24:3.8) {$1$};
  \draw[tgAccent] (n1)--(a)--(b)--(n3);
  \draw[tgAccent] (n1)--(v)--(n3);
\end{tikzpicture}

\small Liczby oznaczają kolory. Wyróżniona krzywa Jordana\\
oddziela sąsiadów o kolorach 2 i 4.
\end{center}

\begin{uwaga}[Twierdzenie o czterech kolorach]
Wynik silniejszy mówi, że każdemu skończonemu grafowi planarnemu
wystarczają cztery kolory. Jego znany dowód korzysta
z obszernej analizy przypadków i weryfikacji komputerowej;
nie przedstawiamy go jako prostego następstwa powyższego
argumentu. Dowód pięciokolorowalności jest natomiast
w pełni elementarny i ilustruje znaczenie twierdzenia Jordana.
\end{uwaga}


% ============================================================
% Iteracja 4, rozdział 11
% ============================================================
